\section{Replace a gcd pair with a smaller pair}
The greatest common divisor of two nonnegative integers is the largest integer that
divides both, apart from the stated zero-pair convention. Division gives
\(A=qB+r\). Any common divisor of \(A\) and \(B\) divides \(r=A-qB\), and any common
divisor of \(B\) and \(r\) divides \(A=qB+r\). Thus replacing \((A,B)\) with
\((B,A\bmod B)\) preserves the set of common divisors. When \(B\) becomes zero,
\(A\) is the answer.

For \((3120,17)\), the successive divisions are
\[
\begin{aligned}
3120&=183\cdot17+9,\\
17&=1\cdot9+8,\\
9&=1\cdot8+1,\\
8&=8\cdot1+0.
\end{aligned}
\]
The pairs at the loop boundary are therefore \((3120,17)\), \((17,9)\), \((9,8)\),
\((8,1)\) and \((1,0)\). Check the second value before taking a remainder: the
last pair is a successful termination, not a request to divide by zero.

\Needspace{14\baselineskip}
\MARSCodeLabel{Replace both values only after saving the remainder.}
\begin{minted}[samepage=true]{mars-idea}
gcd_loop:
    beqz $t1, gcd_done
    nop
    dmodu $t2, $t0, $t1
    move $t0, $t1
    move $t1, $t2
    b gcd_loop
    nop
\end{minted}

Here \texttt{t0} holds \(A\), \texttt{t1} holds \(B\), and \texttt{t2} temporarily
holds the remainder. Writing \texttt{t0} first is safe because the remainder has already
been computed from the old pair. The doubleword unsigned remainder instruction treats
bit 63 as part of the magnitude. A signed remainder would interpret an input at least
\(2^{63}\) as negative. Use \texttt{dmodu} for the full unsigned input contract.

Trace the five pairs above before trying larger values. Then test a zero first
operand, a zero second operand, equal operands and a pair with bit 63 set. Observe
\texttt{v0} as an unsigned bit pattern: a valid gcd can itself have bit 63 set.
